\section{Earlier tests: a brief account of the exclusions}
\label{sec:previous}

The earlier investigations narrowed the required physical readout.
Their exclusions concern the specified constructions, rather than
entire field theories or all possible couplings. The detailed
calculations, sources and diagnostics remain in the cited notes.

\paragraph{WZW evolution and the reflected collar.}
The deterministic SU(2) level-2 pilot gave an explicit selected block
and a positive integrated tensor-norm balance for a constant Loewner
driver. General moving drivers did not preserve contraction, and the
raw smeared kernel did not approach the required identity operator
\cite{PilotNote}. A subsequent full-module collar had a positive norm
balance and a strong identity limit, but its fixed reflected readout
was self-adjoint. A bounded self-adjoint causal scalar operator on an
interval is a multiplication operator, which excludes the nontrivial
arithmetic Volterra response. The selected primary also had the wrong
gamma-tower weights and no first-prime delay \cite{CollarNote}.

\paragraph{Direct arithmetic Loewner sources.}
The completed xi logarithmic derivative generated genuine normalized
Loewner evolutions after a safe spectral translation, with an exact
formula recovering the arithmetic source. The arithmetic information
was supplied through xi itself. Removing the safe translation did not
follow from the geometric construction or its weighted bound, so this
did not furnish an independently motivated physical readout
\cite{SourceNote}.

\paragraph{Brownian bridge and first-passage readouts.}
Bridge-range moments supplied an independent scalar xi identity. The
specified positive first-passage readout was causal and had an ordinary
norm balance, but its Laplace response was not the required quotient:
its kernel had the wrong front behavior and no arithmetic singularity
at $\log2$. A logarithmic readout recovered the boundary phase without
establishing its full-family causality. The distinction was between
the native moment identity and the measured response
\cite{SearchNote,BrownianNote}.

\paragraph{Modular scattering and a cusp load.}
Passing to exact one-forms supplied the rational completion of modular
scattering and realized $K_{1/2}(p)=\xi(p)/\xi(p+1)$ with the stated
ordinary radiation norm and fixed-shift causality. This remains a
positive benchmark. The tested real, frequency-independent cusp delta
load was lossless and causal but had the wrong arithmetic shift
tangent, so it did not extend that benchmark to the required family
\cite{ModularNote}.

\paragraph{Fractional clocks and scalar holonomy.}
A positive weighted diffusion field varied the leading front exponent,
but its scalar clock composition with the modular response failed the
subleading term, boundary modulus, fixed prime delays and identity
endpoint. An analytical exclusion covered scalar complete-Bernstein
clocks, not general coherent couplings \cite{FractionalNote}. A flat
scalar unitary twist on the original modular surface also supplies no
continuous control: $S^2=R^3=1$ allows only six characters, and trivial
cusp holonomy leaves only the trivial character. Nontrivial cusp
holonomy removes the open constant mode \cite{OrbitNote}.

These earlier tests motivated changing primitive orbit weights and the
archimedean response together. Sections~\ref{sec:interface} and
\ref{sec:thermal_boundary} carry that proposal through two explicit
thermal-interface tests. Their failures concern the stated readouts and
fixed-core boundary architecture, while the half-shift benchmark and
thermal coefficient identity remain available.
